\section{Quotients}

\begin{definition}[Strong quasi-projectivity]
\label{akcp-2-1-strong-qproj}
\dcref{2.1}
\source{MR0555258-compactifying-picard:page-0014}
\uses{akcp-1-1-H}
A finitely presented morphism $X\to S$ is strongly quasi-projective (resp.
strongly projective) if $X$ is a retrocompact (resp. closed) subscheme of
$\PP_S(E)$ for a locally free $\OO_S$-module $E$ of constant finite rank.
\end{definition}

\begin{lemma}[Flattening stratification]
\label{akcp-2-3-flattening}
\dcref{2.3}
\source{MR0555258-compactifying-picard:page-0015}
For a finitely presented locally projective $X\to S$, a locally finitely
presented $\FF$ and a rational polynomial $\varphi$, there is a retrocompact
subscheme $Z\subset S$ such that $T\to S$ factors through $Z$ exactly when
$\FF_T$ is $T$-flat with fiber Hilbert polynomial $\varphi$.
\end{lemma}
\begin{proof}
The assertion is local on $S$.  Descend to a noetherian affine approximation
where the standard flattening stratification gives a locally closed $Z_0$ with
the required universal property.  Its inverse image in $S$ is retrocompact and
has the same property by base-change compatibility.
\end{proof}

\begin{lemma}[Regularity in an exact sequence]
\label{akcp-2-4-regularity}
\dcref{2.4}
\source{MR0555258-compactifying-picard:page-0016}
If $0\to\II\to\FF\to\GG\to0$ is exact on a projective scheme and $\FF$ is a
$b$-sheaf, then $\II,\FF,\GG$ are $m$-regular for all $m$ at least a universal
polynomial in $b$ and the coefficients of their Hilbert polynomials.
\end{lemma}
\begin{proof}
The subsheaf $\II$ is a $b$-sheaf, and the main regularity theorem gives a
uniform bound for $\II$ and $\FF$.  The cohomology sequence of
$0\to\II(m-q)\to\FF(m-q)\to\GG(m-q)\to0$ then gives the required vanishing
for $\GG$, uniformly in $m$.
\end{proof}

\begin{definition}[Quotient functors and pseudo-ideal]
\label{akcp-2-5-quot-functor}
\dcref{2.5}
\source{MR0555258-compactifying-picard:page-0016}
For an $S$-flat quotient $\FF\twoheadrightarrow\GG$ with locally finitely
presented target, its pseudo-ideal is $\II(\GG)=\ker(\FF\to\GG)$.  Let
$\Quot_{\FF/X/S}$ assign to $T\to S$ the locally finitely presented,
$T$-flat quotients of $\FF_T$ whose support is proper and finitely presented
over $T$.  The subfunctor $\Quot^{\varphi}_{\FF/X/S}$ imposes fiber Hilbert
polynomial $\varphi$.
\end{definition}

\begin{theorem}[Representability of Quot]
\label{akcp-2-6-quot-representability}
\dcref{2.6}
\source{MR0555258-compactifying-picard:page-0017}
\uses{akcp-2-3-flattening,akcp-2-4-regularity,akcp-2-5-quot-functor}
Suppose $X\to S$ is strongly projective (resp. strongly quasi-projective),
$\FF$ is locally finitely presented, and $\FF$ is a quotient of
$f^*B\otimes\OO_X(v)$ with $B$ locally free of constant rank.  Then
$\Quot^{\varphi}_{\FF/X/S}$ is represented by a universal quotient
$\FF_Q\twoheadrightarrow\GG$ on a strongly projective (resp. strongly
quasi-projective) $S$-scheme $Q$.  For $m$ bounded by a universal polynomial in
the ranks, $v$, and $\varphi$, $f_*\GG(m)$ is locally free of rank $\varphi(m)$
and the Pluecker map embeds $Q$ in
$\PP(\bigwedge^{\varphi(m)}(B\otimes\Sym^{m-v}E))$.
\end{theorem}
\begin{proof}
First embed $X$ as a closed subscheme of $\PP(E)$ and replace $\FF$ by
$f^*B\otimes\OO(v)$.  A quotient of this sheaf on $\PP(E)$ restricts to a
quotient on $X$ precisely when the induced map from the kernel of
$f^*B\otimes\OO(v)\to\FF$ to the target vanishes.  By
\cref{akcp-1-1-H}, vanishing is cut out by a closed subscheme, so the desired
functor is a closed subfunctor of the ambient Quot functor.

Choose $m$ using \cref{akcp-2-4-regularity}; every quotient and pseudo-ideal
is then $m$-regular on fibers.  The universal quotient therefore gives a map
to the Grassmannian of $\varphi(m)$-quotients of $f_*\FF(m)$.  Its Pluecker
coordinates satisfy the usual relations, and the Hilbert polynomial condition
is closed by \cref{akcp-2-3-flattening}.  The resulting closed subscheme
represents the Quot functor; the displayed embedding and formula for
$\OO_Q(1)$ follow from the Pluecker construction.  The valuative criterion for
properness gives strong projectivity in the projective case.

For quasi-projective $X$, take its scheme-theoretic closure $\bar X$ in
$\PP(E)$.  Extend $\FF$ by zero and take the locally finitely presented image
on $\bar X$.  The ambient Quot scheme above contains the desired functor as
the open locus where the universal support avoids $\bar X\setminus X$; this
locus is open because the projection is proper.  Hence it is strongly
quasi-projective.
\end{proof}

\begin{corollary}[Locally projective Quot]
\label{akcp-2-7-local-quot}
\dcref{2.7}
\source{MR0555258-compactifying-picard:page-0021}
\uses{akcp-2-6-quot-representability}
For finitely presented locally projective (resp. locally quasi-projective)
$X\to S$, $\Quot_{\FF/X/S}$ is a disjoint union of locally finitely presented,
locally projective (resp. locally quasi-projective) $S$-schemes.
\end{corollary}
\begin{proof}
This is local on $S$; on an affine open, apply the theorem to each Hilbert
polynomial and use the standard open-and-closed decomposition by polynomial.
\end{proof}

\begin{corollary}[Hilbert schemes]
\label{akcp-2-8-hilbert}
\dcref{2.8}
\source{MR0555258-compactifying-picard:page-0021}
\uses{akcp-2-6-quot-representability}
For strongly projective (resp. strongly quasi-projective) $X/S$, the functor
$\Hilb^\varphi_{X/S}$ is represented by a strongly projective (resp. strongly
quasi-projective) $S$-scheme.
\end{corollary}
\begin{proof}
Apply \cref{akcp-2-6-quot-representability} to $\FF=\OO_X$ and the zero twist.
\end{proof}

\begin{theorem}[Quotient by a proper equivalence relation]
\label{akcp-2-9-effective-quotient}
\dcref{2.9}
\source{MR0555258-compactifying-picard:page-0021}
\uses{akcp-2-8-hilbert}
Let $X/S$ be strongly quasi-projective and let $R\subset X\times_S X$ be a
flat, finitely presented, proper equivalence relation.  Assume the projections
$R\to X$ have only finitely many Hilbert polynomials.  Then $R$ is effective;
$q:X\to X/R$ is faithfully flat and strongly projective; and $X/R\to S$ is
strongly quasi-projective.
\end{theorem}
\begin{proof}
Let $H$ be the disjoint union of the finitely many Hilbert schemes supplied by
\cref{akcp-2-8-hilbert}, with universal subscheme $W\subset X\times_S H$.
The relation $R$ determines a unique map $g:X\to H$ for which
$(1\times g)^*W=R$.  For two $T$-points $x_1,x_2$ of $X$,
\[
x_1\,R\,x_2\quad\Longleftrightarrow\quad g(x_1)=g(x_2),
\]
because equality of the corresponding fibers of $W$ is equivalent to mutual
membership in the reflexive, symmetric, transitive relation.

The graph $\Gamma_g\subset X\times_S H$ is closed (since $H/S$ is separated),
and Step III of the source identifies it with a finitely presented closed
subscheme of $W$.  The projection $W\to H$ is flat, quasi-compact and
surjective; descent along this faithfully flat map identifies the two pullbacks
of $\Gamma_g$, so faithfully flat descent supplies a finitely presented
subscheme $Z\subset H$.  The map $g$ factors through $Z$, and
$X\times_ZX=R$.

The map $X\to Z$ is faithfully flat, finitely presented and strongly projective
because it is the pullback of $W\to H$.  It is a universal effective epimorphism,
so $Z$ is the quotient by $R$ and the asserted properties follow.
\end{proof}

\begin{corollary}[Locally projective quotient]
\label{akcp-2-10-local-quotient}
\dcref{2.10}
\source{MR0555258-compactifying-picard:page-0023}
\uses{akcp-2-9-effective-quotient}
If $X/S$ is locally projective and $R$ is flat, finitely presented, proper and
an equivalence relation, then $R$ is effective, $X\to X/R$ is faithfully flat,
finitely presented and proper, and the quotient is locally quasi-projective.
\end{corollary}
\begin{proof}
Work over an affine open of $S$, where $X/S$ is strongly quasi-projective, and
apply \cref{akcp-2-9-effective-quotient}; the properties descend on the base.
\end{proof}
